\scanpage{85}


est inversible, alors le noyau de $*:E_1\otimes E_2\longrightarrow E_1*E_2$ est $\operatorname{Im}({\rm id}_{E_1\otimes E_2}-\sigma_1\otimes\sigma_2)$. [\uline{NB} Si $T_1$ et $T_2$ sont nuls, on a $\sigma_i={\rm id}$ $(i=1,2)$ donc $\mathrm{id}-\sigma_1\otimes\sigma_2=0$,
dont l'image est nulle, alors que le noyau de $E_1\otimes E_2\longrightarrow E_1*E_2$ \deleted{\unsure{contient}} \inserted{est ég. à} \inserted{celui de l'augmentation,} $E_1*_{0}E_2=L_1\otimes E_2+E_1\otimes L_2$ (qui est de rang $3$ si si $L_1,L_2$ inv.).]

En termes des scindages $x_1,x_2$ de $E_1,E_2$, et des \deleted{\unsure{la}} décompositions correspondantes

\[ \begin{aligned} E_1&\simeq L_1\oplus O\,x_1,\quad E_2\simeq L_2\oplus\underline O\,x_2\\ E_1\otimes E_2&\simeq(L_1\otimes L_2)\oplus\deleted{\unsure{\underline O}}\underset{\simeq L_1}{\underline{L_1\otimes x_2}}\oplus\underset{\simeq L_2}{\underline{x_1\otimes L_2}}+\underset{\simeq O}{\underline{\underline O\cdot x_1\otimes x_2}} \end{aligned} \]

et compte tenu que $\sigma_1\otimes\sigma_2\inserted{-{\rm id}}$ est \deleted{l'id.} \inserted{nul} sur le morceau $L_1\otimes L_2$, \deleted{il reste à voir} $\operatorname{Im}(\deleted{\unsure{{\rm id}-}}\sigma_1\otimes\sigma_2-{\rm id}_{E_1\otimes E_2})$ est engendré par $3$ sortes d'éléments, correspondants aux trois autres facteurs

\[ \begin{aligned} \overline{u_1\otimes x_2}-u_1\otimes x_2&=-u_1\otimes\underset{\substack{\Vert\\T_2-x_2\\\Vert\\x_2-b_2}}{\overline{x_2}}\unsure{-|=}u_1\otimes x_2=-\underline{2u_1}\otimes x_2+\underline{u_1\otimes b_2}\\ \overline{x_1\otimes u_2}-x_1\otimes u_2&=-2x_1\otimes\underline{u_2}+b_1\otimes u_2\\ \overline{x_1\otimes x_2}-x_1\otimes x_2&=\underset{(x_1-b_1)\otimes(x_2-b_2)}{\overline{x_1}\otimes\overline{x_2}}-x_1\otimes x_2=-x_1\otimes b_2-b_1\otimes x_2+b_1\otimes b_2 \end{aligned} \]

donc les éléments (avec les notations de la somme directe)

\[ \left\{\begin{aligned}&-2u_1+u_1\otimes b_2 &&u_1\in\unsure{T|\Gamma} L_1\\ &-2u_2+b_1\otimes u_2 &&u_2\in\Gamma L_2\\ &-b_1-b_2+b_1\otimes b_2\end{aligned}\right. \]

En car. $2$ cela se réduit à

\[ \left\{\begin{aligned}&u_1\otimes b_2,\ b_1\otimes u_2\in\Gamma L_1\otimes L_2\\ &-b_1-b_2-b_1\otimes b_2\quad\unsure{\notin}\end{aligned}\right. \]
% archive sheet number: 84
